\section*{第 \ref{ch:g12:rc-rl-circuits} 章——RC 与 RL 电路}

\begin{solution}{exo:g12:rc-rl-circuits:1}
$q = Cu = 4.7\times10^{-4} \times 12 \approx \qty{5.6}{mC}$;
$u = q/C = 2.0\times10^{-3}/10^{-4} = \qty{20}{V}$.
\end{solution}

\begin{solution}{exo:g12:rc-rl-circuits:2}
\qty{47}{pF}：收音机调谐器；\qty{2200}{\micro F}：平滑用；\qty{10}{F}：
超级\omterm{def:g12:rc-rl-circuits:capacitor}{电容器}。$E = \tfrac12 \times 10 \times 2.7^2 \approx \qty{36}{J}$。
\end{solution}

\begin{solution}{exo:g12:rc-rl-circuits:3}
$i = C\,\Delta u/\Delta t = 2.2\times10^{-4} \times 5.0/0.010 =
\qty{0.11}{A}$。然后 $i = C\,du/dt = 0$：带电\omterm{def:g12:rc-rl-circuits:capacitor}{电容器}相当于
断开的电路——无稳定电流。
\end{solution}

\begin{solution}{exo:g12:rc-rl-circuits:4}
$RC = 10^4 \times 10^{-4} = \qty{1.0}{s}$;
$L/R = 0.10/50 = \qty{2.0}{ms}$. $\unit{\ohm}\,\unit{F} =
(\unit{V/A})(\unit{C/V}) = \unit{C/A} = \unit{s}$;
$\unit{H}/\unit{\ohm} = (\unit{V.s/A})/(\unit{V/A}) = \unit{s}$.
\end{solution}

\begin{solution}{exo:g12:rc-rl-circuits:5}
$3.8/6.0 = 0.63$，故 $\tau = \qty{2.0}{s}$；$95\%$ 对应 $3\tau =
\qty{6.0}{s}$.
\end{solution}

\begin{solution}{exo:g12:rc-rl-circuits:6}
$\tau = 4700 \times 2.2\times10^{-4} \approx \qty{1.0}{s}$;
$u(\tau) = 0.63 \times 9.0 = \qty{5.7}{V}$;
$u(3\tau) = 0.95 \times 9.0 = \qty{8.6}{V}$;
$I_0 = \mathcal{E}/R = 9.0/4700 \approx \qty{1.9}{mA}$（\omterm{def:g12:rc-rl-circuits:capacitor}{电容器}未充电时：
$\mathcal{E}$ 全部加在 $R$ 两端）。
\end{solution}

\begin{solution}{exo:g12:rc-rl-circuits:7}
$du/dt = (\mathcal{E}/RC)e^{-t/RC}$，故 $RC\,du/dt + u =
\mathcal{E}e^{-t/RC} + \mathcal{E}(1 - e^{-t/RC}) = \mathcal{E}$.
$u(0) = 0$：\omterm{def:g12:rc-rl-circuits:capacitor}{电容器}开始时不带电；$u \to \mathcal{E}$：
充满电，电流已停止。
\end{solution}

\begin{solution}{exo:g12:rc-rl-circuits:8}
$\tau = 4.7\times10^{4} \times 10^{-4} = \qty{4.7}{s}$;
$u(\tau) = 12\,e^{-1} \approx \qty{4.4}{V}$；\qty{0.60}{V} 为 $5\%$，
在 $3\tau \approx \qty{14}{s}$ 时达到。
\end{solution}

\begin{solution}{exo:g12:rc-rl-circuits:9}
$u = L\,\Delta i/\Delta t = 0.50 \times 2.0/0.040 = \qty{25}{V}$.
闭合时，$i$ 按电路本身的快慢增长，$di/dt$ 较小；断开时
\emph{迫使} $i$ 几乎瞬间降为零，$di/dt$ 极大，$u$
达到足以产生火花的数值。
\end{solution}

\begin{solution}{exo:g12:rc-rl-circuits:10}
$di/dt = (\mathcal{E}/L)e^{-Rt/L}$，故 $L\,di/dt + Ri =
\mathcal{E}e^{-Rt/L} + \mathcal{E}(1 - e^{-Rt/L}) = \mathcal{E}$.
$i(\tau) = 0.63 \times 2.0 = \qty{1.3}{A}$;
$i(3\tau) = 0.95 \times 2.0 = \qty{1.9}{A}$.
\end{solution}

\begin{solution}{exo:g12:rc-rl-circuits:11}
$u'(t) = (\mathcal{E}/\tau)e^{-t/\tau}$，故 $u'(0) = \mathcal{E}/\tau$：
切线 $u = (\mathcal{E}/\tau)\,t$ 到达 $\mathcal{E}$ 的时刻为
$t = \tau$. $C = \tau/R = 0.50/2000 = \qty{250}{\micro F}$.
\end{solution}

\begin{solution}{exo:g12:rc-rl-circuits:12}
超级\omterm{def:g12:rc-rl-circuits:capacitor}{电容器} $\tfrac12 \times 3000 \times 2.7^2 \approx \qty{1.1e4}{J}$；
闪光灯 $\tfrac12 \times 8.0\times10^{-4} \times 300^2 = \qty{36}{J}$；
\omterm{def:g11:magnetic-fields:solenoid}{线圈} $\tfrac12 \times 0.010 \times 10^2 = \qty{0.50}{J}$。电池
（$\qty{1.3e4}{J}$）$>$ 超级\omterm{def:g12:rc-rl-circuits:capacitor}{电容器} $\gg$ 闪光灯 $\gg$ \omterm{def:g11:magnetic-fields:solenoid}{线圈}——但三者
都能在毫秒内释放\omterm{def:g10:energy-conservation:energy}{能量}：关心的是\omterm{def:g11:work-of-force:power}{功率}，而非\omterm{def:g10:energy-conservation:energy}{能量}。
\end{solution}

\begin{solution}{exo:g12:rc-rl-circuits:13}
$\tau = \qty{5.0}{s}$，故 $R = \tau/C = 5.0/10^{-4} = \qty{50}{k\ohm}$。
$1 - e^{-t/\tau} = 0.50$ 得 $e^{-t/\tau} = 0.50$，$t = \tau\ln 2
\approx 0.69\,\tau = \qty{3.5}{s}$.
\end{solution}

\begin{solution}{exo:g12:rc-rl-circuits:14}
$\Delta u = i\,\Delta t/C = 0.50 \times 0.010/2.2\times10^{-3} \approx
\qty{2.3}{V}$。需满足 $C \geq i\,\Delta t/\Delta u =
5.0\times10^{-3}/0.50 = \qty{10}{mF}$.
\end{solution}

\begin{solution}{exo:g12:rc-rl-circuits:15}
$u \approx L\,\Delta i/\Delta t = 0.80 \times 2.0/10^{-3} =
\qty{1.6}{kV}$; $E_L = \tfrac12 \times 0.80 \times 2.0^2 =
\qty{1.6}{J}$。缓慢断开时，$di/dt$ 保持较小，$u$ 仅为几伏：
不会产生火花。
\end{solution}

\begin{solution}{pb:g12:rc-rl-circuits:1}
\textbf{1.} $q = CU_0 = 8.0\times10^{-4} \times 300 = \qty{0.24}{C}$.

\textbf{2.} $N = 0.24/1.60\times10^{-19} = \num{1.5e18}$.

\textbf{3.} $E_C = \tfrac12 \times 8.0\times10^{-4} \times 300^2 =
\qty{36}{J}$.

\textbf{4.} $\num{1.3e4}/36 \approx 360$ 倍。

\textbf{5.} 电池储存的\omterm{def:g10:energy-conservation:energy}{能量}很多，但释放缓慢；\omterm{def:g12:rc-rl-circuits:capacitor}{电容器}
把不多的\omterm{def:g10:energy-conservation:transfer}{能量转化}为极大的\emph{\omterm{def:g11:work-of-force:power}{功率}}。

\textbf{6.} $\mathcal{E} = Ri + u$，其中 $i = C\,du/dt$：
$RC\,du/dt + u = \mathcal{E}$.

\textbf{7.} $du/dt = (\mathcal{E}/RC)e^{-t/RC}$：
$\mathcal{E}e^{-t/RC} + \mathcal{E}(1 - e^{-t/RC}) = \mathcal{E}$
对所有 $t$ 都成立，且 $u(0) = 0$。

\textbf{8.} $\tau = RC = 10^3 \times 8.0\times10^{-4} = \qty{0.80}{s}$。

\textbf{9.} 当 $3\tau = \qty{2.4}{s}$ 时，$u = 0.95 \times 300 =
\qty{285}{V}$——两秒的嗡嗡声。

\textbf{10.} $I_0 = \mathcal{E}/R = \qty{0.30}{A}$；随 $u$ 升高，
余下的 $\mathcal{E} - u$ 加在 $R$ 上，随之减小，故 $i = (\mathcal{E}-u)/R$
逐渐\omterm{def:g12:mechanical-waves:attenuation}{衰减}至零。

\textbf{11.} $u(t) = U_0\,e^{-t/\tau'}$，其中 $\tau' = R_f C =
0.50 \times 8.0\times10^{-4} = \qty{0.40}{ms}$。

\textbf{12.} $I_0 = U_0/R_f = 300/0.50 = \qty{600}{A}$。

\textbf{13.} $P_0 = U_0 I_0 = 300 \times 600 = \qty{180}{kW}$——约为
$80$ 只电热水壶同时工作的\omterm{def:g11:work-of-force:power}{功率}。

\textbf{14.} 时长 $3\tau' = \qty{1.2}{ms}$；平均\omterm{def:g11:work-of-force:power}{功率}
$\approx 36/1.2\times10^{-3} = \qty{30}{kW}$。

\textbf{15.} $2.4/1.2\times10^{-3} = 2000$。仅\omterm{def:g11:circuits-and-power:resistance}{电阻}发生了变化
（接入 \qty{1.0}{k\ohm}，改为 \qty{0.50}{\ohm}）：一个\omterm{def:g11:work-of-force:power}{功率}杠杆。

\textbf{16.} $C = 2E_C/U_0^2 = 400/1500^2 \approx \qty{1.8e-4}{F} =
\qty{180}{\micro F}$.

\textbf{17.} $q = CU_0 = 1.78\times10^{-4} \times 1500 \approx
\qty{0.27}{C}$.

\textbf{18.} $\tau = 50 \times 1.78\times10^{-4} \approx \qty{8.9}{ms}$,
时长 $3\tau \approx \qty{27}{ms}$——约为闪光时长的二十倍。

\textbf{19.} $I_0 = 1500/50 = \qty{30}{A}$;
$P_0 = 1500 \times 30 = \qty{45}{kW}$.

\textbf{20.} 由 $3\tau_c = \qty{6.0}{s}$ 得 $\tau_c = \qty{2.0}{s}$，故
$R_c = \tau_c/C = 2.0/1.78\times10^{-4} \approx \qty{11}{k\ohm}$，且
$I_0 = 1500/1.12\times10^{4} \approx \qty{0.13}{A}$——六秒即可
充好电，只用十分之一\omterm{def:g11:circuits-and-power:current}{安培}的电流。
\end{solution}
